\scanpage{79}
% manuscript line 1
applique la prop à $f'$, et on \unsure{applique ens.|appliquera|applique} \unsure{le $=$|$=$} à $f''$.

% manuscript line 2
\deleted{Prop} \uline{Corollaire 2} Sous les conditions de la prop. pour $\underline W$, \deleted{$\unsure{\mathfrak X|\mathscr X}$ \unsure{Si}} supposons que \deleted{$\mathfrak X$ et $\unsure{\mathfrak Y|\mathscr Y}$} tte partie finie de $\mathfrak X$ soit contenue dans un ouvert affine. Alors si $f:\mathfrak X\to\mathfrak Y$ est séparé, alors il en est de même de $\underline W^{\unsure{\bullet|a}}(f)$ [sous réserve que $\underline W^\wedge(\mathfrak Y)$ existe dans le cas $\mathrm I'$].

% manuscript line 9
On applique la prop à $\mathfrak X\unsure{\text{---}|\longrightarrow}\mathfrak X\times_{\mathfrak Y}\mathfrak X$.

% manuscript line 10
\uline{Proposition} $W$, \deleted{$\unsure{\mathcal U|\mathfrak X}$} $\mathfrak Y$ comme dans la prop. précédente. Soit $f:\mathfrak X\to\mathfrak Y$, on suppose \deleted{\unsure{déjà|affine}} \inserted{dans le cas $\mathrm I'$} que pour tout ouvert affine \inserted{$U$} de $\mathfrak Y$, \inserted{tte partie finie de} $f^{-1}(U)$ est \deleted{réunion} contenue dans un ouvert affine (p. ex. $f$ quasi-projectif). On suppose de plus $\mathfrak Y$ loc. noeth et \deleted{$\mathfrak X$.} $f$ de type fini. \inserted{\deleted{(resp. affine)}} Alors $\underline W^{a}(f)$ est de type \unsure{finie|fini} \deleted{(resp. affine)}.
% slanted
\margin{\unsure{rédiger cela|restreig. à cela}. \unsure{Donner|Dans} le cas ``Greenberg spécial''\unsure{;|,} l'autre est \unsure{une} variante \unsure{contenue|continue|remplacée} dans le premier \unsure{(car} projectif $\Rightarrow$ quasi-projectif\unsure{)}}.
% archive sheet number: 78
